\section{观测器设计}

未知方向的控制中的一个难点在于 $B_i$ 未知，会导致无法获取到自身的加速度。
而在利用部分相对位置观测相对位置和速度的观测器中，自身的加速度需要部分补充道观测器中，所以无法使用现有的常值参数的观测器。
考虑到现在的加速度为$\ddot{p}_i=B_i u_i + d_i$， 我们只能将整个加速度当作扰动来处理，其上界为：$\|\ddot{p}_i\| \leq \bar{b} \|u_i\| + \bar{d}$，事情的难点就在于$\|u_i\|$ 中的幅值 $R_i(t)$ 包含势场项目$\phi_i$。
当作扰动设计观测器的话，观测器就必须有一个参数需要时变并且大于$\phi_i$。
这会给通信连通的证明带来许多困难。

大体算法设计，请参考


大体证明思路。

1. 假设存在一个时刻通信断开了。
2. 在这个时刻前观测器是收敛的，所以观测误差是有界的。
3. 有届的观测误差下，系统不会跑飞，也就是说这个时刻之前应该没有断开
4. 根据位置的连续性可知，不可能断开连通性。